% Generated draft: requires mathematical and source review.
\documentclass[11pt]{article}
\usepackage[margin=25mm]{geometry}
\usepackage{fontspec}
\setmainfont{Latin Modern Roman}
\usepackage{amsmath,amssymb,mathtools,unicode-math}
\setmathfont{Latin Modern Math}
\usepackage{enumitem,needspace,xurl,hyperref}
\setlength{\parindent}{0pt}
\setlength{\parskip}{5pt}
\setlength{\emergencystretch}{3em}
\newcommand{\N}{\mathbb{N}}
\newcommand{\Z}{\mathbb{Z}}
\newcommand{\Q}{\mathbb{Q}}
\newcommand{\R}{\mathbb{R}}
\newcommand{\C}{\mathbb{C}}
\newcommand{\F}{\mathbb{F}}
\newcommand{\A}{\mathbb{A}}
\newcommand{\PP}{\mathbb{P}}
\newcommand{\E}{\mathbb{E}}
\newcommand{\eps}{\varepsilon}
\newcommand{\dd}{\,\mathrm{d}}
\begin{document}
\section*{UnsolvedMath ID: 30003270}
\textbf{Source problem code:} catalogue entry
\subsection{Definitions and mathematical statement}
The supplied cleaned formulation uses a positive square-free integer \(d\) and the imaginary quadratic field \(L=\Q(\sqrt{-d})\). Write \(\overline{x}\) for its nontrivial \(\Q\)-automorphism and \(g^*=\overline{g}^{\,t}\) for conjugate transpose. The record specifies
\[
G(\Q)=\{g\in\operatorname{SL}_3(L):g^*Jg=J\}
\]
and assumes that the corresponding algebraic group \(G\) has \(\Q\)-rank one, meaning that a maximal \(\Q\)-split torus in \(G\) has dimension one. For the displayed equation to define the asserted unitary group, the conventional interpretation requires \(J\in\operatorname{GL}_3(L)\) with \(J^*=J\). The record does not specify \(J\); these conditions explain the unitary interpretation and do not identify an intended matrix.

Let \(\Gamma\leq G(\Q)\) be an arithmetic subgroup, as required by the cleaned formulation. Here arithmetic means that, in a faithful algebraic representation \(\rho:G\hookrightarrow\operatorname{GL}_N\) defined over \(\Q\), the subgroup \(\rho(\Gamma)\) is commensurable with \(\rho(G(\Q))\cap\operatorname{GL}_N(\Z)\); two subgroups are commensurable when their intersection has finite index in each.

\emph{Source question, with unresolved qualifications.} For groups and arithmetic subgroups satisfying these assumptions, can one carry out explicit reduction theory for \(\Gamma\)? The record does not define the required output of an explicit reduction theory, specify which choices of \(J\) are intended, or say whether one uniform procedure for all such inputs is required. Consequently, the question cannot yet be expressed as a uniquely specified existence or algorithmic assertion.

There is also a substantive discrepancy between the supplied formulations: the original statement uses \(L=\Q(\sqrt d)\), with \(d>0\) square-free, and merely calls \(\Gamma\) a subgroup, whereas the cleaned statement uses \(L=\Q(\sqrt{-d})\) and requires \(\Gamma\) to be arithmetic. The original also describes \(G\) as a \(\Q\)-rank-one form of \(\operatorname{SL}_3\). The imaginary quadratic and arithmetic assumptions above reproduce the cleaned formulation; the supplied evidence does not resolve their compatibility with the original.
\subsection{Short English statement}
The cleaned question asks for explicit reduction theory for an arithmetic subgroup of a \(\Q\)-rank-one unitary group in three variables over an imaginary quadratic field. The Hermitian form and the intended reduction outputs remain unspecified. The original wording instead uses a positive radicand and does not require the subgroup to be arithmetic; these differences require source review.
\subsection{Sources}
\begin{itemize}
\item Local UnsolvedMath registry, supplied complete record, permanent ID \(30003270\), source problem code \texttt{catalogue entry}, title: \emph{Explicit Reduction Theory for Rank-One Unitary Arithmetic Groups}. Collection: Oberwolfach Reports Open Problems. No individual proposer is recorded; the supplied proposed year is 2016.
\item Supplied report citation: \emph{Mini-Workshop: Computations in the Cohomology of Arithmetic Groups (2017)}, Oberwolfach Reports, pp.\ 2941--2973, DOI 10.4171/owr/2016/52. Source URL and supplied literature source ``OWR 52 report'': \url{https://doi.org/10.4171/owr/2016/52}. The citation's 2017 label and the record's 2016 year information are retained without reconciliation.
\item Supplied literature source ``Yasaki related reduction'': \url{https://arxiv.org/abs/math/0601071}. No further bibliographic details are supplied.
\item Literature assessment dated exactly \texttt{2026-08-21}: the supplied curation labels the problem open and reports that explicit reduction is known in some related unitary examples, but that no general explicit reduction theory for the groups posed in the report was found. This is supplied dated curation, not an independently verified literature assessment. The background explicitly requests independent verification of time-sensitive status and cited claims.
\item The registry labels the statement \texttt{corrected\_verified} and says that mathematical content was preserved while notation, typography, or omitted definitions were repaired. This is supplied editorial metadata, not independent verification of the differing formulations or mathematical correctness.
\end{itemize}
\end{document}
